\setcounter{Exercise}{0}

\chapter*{Feuille d'exercices : Les fonctions}

\begin{Exercise}[title={Appréhendez une définition de fonction}]

\begin{lstlisting}[style=pythonstyle, language=Python]
def compteur3():
    i = 0
    while i < 3:
        print(i)
        i = i + 1

def double_compteur3():
    compteur3()
    compteur3()

def compteur(stop):
    i = 0
    while i < stop:
        print(i)
        i = i + 1

def compteur_complet(start, stop, step):
    i = start
    while i < stop:
        print(i)
        i = i + step

def calcule_cube(a):
    cube = a * a * a # ou bien a**3
    return cube

def somme_produit(a,b):
    """Calcule la somme et le produit de deux nombres"""
    s = a + b
    p = a * b
    return s, p
\end{lstlisting}
\Question Pour chacune des fonctions, donner le nom de la fonction, le nombre et le nom des paramètres, le nom des variables locales et le nombre de valeurs de retour.
\Question Pour chacune des fonctions, donner un exemple d'appel vous stockerez les éventuelles valeurs de retour dans des variables.
\end{Exercise}



\newpage

\begin{Exercise}[title={Il y a un truc qui cloche}]
\Question Pour chacun des codes suivants, découvrait ce qui cloche dans la définition ou dans l'appel de la fonction.
\begin{lstlisting}[style=pythonstyle, language=Python]
def aire_carre(longueur):
    return longueur**2

aire = aire_carre()
print("L'aire du carre est de", aire, "m²")
\end{lstlisting}

\begin{lstlisting}[style=pythonstyle, language=Python]
def perimetre_carre(longueur):
    print(4*longueur)

perimetre = perimetre_carre(12)
print("Le perimetre du carre est", perimetre, "m")
\end{lstlisting}

\begin{lstlisting}[style=pythonstyle, language=Python]
def quotient(x,y)
    return x/y 
    
print(quotient(1.7,2.5))
\end{lstlisting}

\begin{lstlisting}[style=pythonstyle, language=Python]
def aire_triangle_rectangle(base,hauteur):
    return base*hauteur/2 

aire_triangle_rectangle(1.7,2.5)
\end{lstlisting}

\begin{lstlisting}[style=pythonstyle, language=Python]
def distance_euclidienne(x1, x2, y1, y2):
    dx = (x1-x2)
    dy = (y1-y2)
    distance = math.sqrt(dx**2 + dy**2)

x1, y1, x2, y2 = 1, 2, 5, 7
distance = distance_euclidienne(x1, x2, y1, y2)
print(f"La distance euclidienne entre ({x1}, {y1}) et ({x2}, {y2}) est {distance} m")
\end{lstlisting}

\begin{lstlisting}[style=pythonstyle, language=Python]
import random

def de():
    tirage = random.randrange(1, 6)
    return tirage

print("Les trois des sont :", DE(), DE(), "et", DE())
\end{lstlisting}

\begin{lstlisting}[style=pythonstyle, language=Python]
def volume_pave(l, l, h):
    return l * l * h

print(volume_pave(2, 72, 10))
\end{lstlisting}

\begin{lstlisting}[style=pythonstyle, language=Python]
def n_print(s, n):
    for i in range(n):
        print(s)

print(volume_pave(2, 72, 10))
\end{lstlisting}

\end{Exercise}

\vfill

\begin{Exercise}[title={Vous n'auriez pas vu la sortie }]
\Question Pour chacun des codes ci-dessous, dire ce qui sera écrit dans le terminal à la fin de l'exécution du code.
\begin{lstlisting}[style=pythonstyle, language=Python]
def n_print(s, n):
    for i in range(n):
        print(s)

n_print("Coucou les loulous", 12)
\end{lstlisting}

\begin{lstlisting}[style=pythonstyle, language=Python]
def est_premier(num):
    if num > 2 and num % 2 == 0:
        return False
    for x in range(2, num // 2):
        if num % x == 0:
            return False
    return True 

def premiers(x):
    for i in range(2, x+1):
        if est_premier(i):
            print(i)

premiers(22)
\end{lstlisting}

\begin{lstlisting}[style=pythonstyle, language=Python]
def f(x, y, z):
    return x + y, x + z, y + z

def g(x, y):
    if x > 10 :
        return 1/y, x/y
    else :
        return 1/x, y/x

def combo(x, y, z):
    a, b, c = f(x, y, z)
    d, e = g(z, x)
    u, v, w = f(a, b, e**2)
    return v - w, b - a 

print(f"combo(4, 9, 2) = {combo(4 ,9 ,2)}")
\end{lstlisting}

\begin{lstlisting}[style=pythonstyle, language=Python]
def triangle(c, x):
    for i in range(1, x):
        print(c * i)
    for i in range(x, 1, -1):
        print(c * i)

def pixel_art(n):
    for i in range(1, n, 2):
        triangle("=", i)
    
pixel_art(7)
\end{lstlisting}
\end{Exercise}



\begin{Exercise}[title={Les fonctions à trous }]
\Question Compléter les codes ci-dessous.
\begin{lstlisting}[style=pythonstyle, language=Python]
def pair_entre_deux_nombres(...):
    for i in range(a, b):
        if ...:
            print(...)

def factorielle(x): # 1 * 2 * ... * x-1 * x
    produit = ...
    for i in range(..., ...):
        produit = produit * ...
    return

def inverser(x, y):
    temp = x
    y = ...
    x = ...
    return ...

def maximum(a, b, c):
    if a > b :
        ...
            return c
        else :
            return ...
    else :
        if c > b :
            return ...
        ...
            return b
\end{lstlisting}

\end{Exercise}

\begin{Exercise}[title={Ça va fonctionner!}]
\Question Coder une fonction somme qui retourne la somme de 2 nombres.
Exemple d'appel:\\
$>>>$ somme(2, 3)\\
5
\Question Coder une fonction plus\_grand qui retourne la plus grande valeur de deux valeurs entières données en paramètres.

Exemple d'appel:\\
$>>>$ plus\_grand(2, 3)\\
3

\Question Coder une fonction repete qui prend deux paramètres mot et n et qui affiche le même mot n fois.

Exemple d'appel:\\
$>>>$ repete("Tartuffe", 3)\\
Tartuffe\\
Tartuffe\\
Tartuffe

\end{Exercise}

% \begin{Exercise}[title={Du global dans mon local}]
% \Question Pour chacun des codes ci-dessous, dire ce qui sera écrit dans le terminal à la fin de l'exécution du code.
% \begin{lstlisting}[style=pythonstyle, language=Python]
% a = 7
% def demi_carre(a):
%     print(a)
%     a = a / 2
%     print(a ** 2)
% print(a)
% print(demi_carre(a))
% \end{lstlisting}

% \begin{lstlisting}[style=pythonstyle, language=Python]
% def f(a, b):
%     b = 12
%     print(b)
% b = 10
% print(f(5, b))
% print(b)
% \end{lstlisting}
% \end{Exercise}